\section{退学潮：不是反学校，而是 applied research 被迫面对市场}

退学是这期节目最显眼的话题，但赵晨阳和主持人都没有把它包装成青春叛逆故事。赵晨阳说自己读了十五个月博士，这十五个月与做 SGLang 的时间高度重合。后来觉得 timing 正确，工业界机会比学术界更多，于是选择离开。这个决定背后的关键词不是“退学”，而是“什么样的 research 应该在哪里发生”。

他把 research 分成 applied research 和 fundamental research。Applied research 既然叫 applied，就必须有 application，必须真正接受市场、用户、工程系统和商业价值的检验。Fundamental research 则更适合学校，因为它可以暂时不在乎市场结果，而追求内在哲学体系的自洽和长期问题意识。问题在于 LLM 时代资源门槛抬高，即使是 fundamental research，要验证很多假设也需要大量算力和工程资源，学界越来越难提供。

\begin{importantbox}{赵晨阳对 applied research 的判断}
Applied research 最大的点就是要有 application。若一个系统方向最终要靠大厂 adoption、真实用户和商业效果来证明价值，那么在学校里虚空 formulate 一个 evaluation，可能反而不如直接到工业界接受真实反馈。
\end{importantbox}

这里还有一个地域差异。赵晨阳认为，美国读博的中国学生会更挣扎，因为签证和工作许可限制让他们很难在读博期间进入最前沿公司，接触最核心技术。相比之下，国内很多模型公司愿意给非正式编制的实习生更高权限；香港、新加坡或国内读博也不一定有同样强烈的离开必要。退学潮因此不是一种普遍建议，而是某些地区、某些身份、某些技术窗口叠加后的结果。

主持人补充了“离经叛道”的文化感：在中国传统观念里，退学常被视为不光彩；而硅谷创业叙事里，退学反而可能被理解成 outlier 特质。赵晨阳也承认，愿意退学的人往往更不循规蹈矩，这种特质对创业和投资人都有吸引力。但他并没有把退学神化。很多人留在学校也完全合理，关键在于你相信的问题、资源条件和机会窗口是否仍然匹配学校环境。

\begin{warningbox}{不要把退学潮误读成“读博无用”}
这期节目真正讨论的不是“该不该退学”，而是不同类型 research 的最优承载环境。Fundamental research 仍然需要学校，系统型 applied research 若必须经受真实 adoption 和商业反馈，则可能更适合工业界。
\end{warningbox}

尚晏仪引用戴雨森的意思，说创业需要某种情绪和信念。如果你非常冷静地推演终局，可能永远不会出来创业，因为总能想到大厂会吃掉你、上下游会挤压你、scope 会越收越小。创业不是完全理性规划出来的，它需要某种相信、冲动和愿意承担不确定性的能力。这也解释了为什么年轻人反而可能更敢跳出来：他们没有被长周期体制完全固定，还保留对 1\% 机会的偏好。

\subsection{本章小结}

退学潮不是反学校，而是 AI 时代 applied research 与工业资源、真实反馈、市场验证之间的关系变得更直接。赵晨阳的选择说明，若一个方向本质上要通过工程 adoption 和商业价值证明，那么工业界可能比学校更本质；但 fundamental research、长期理论问题和自洽体系仍然有学校价值。退学不是模板，而是时机、身份、资源与问题类型共同作用的结果。

